\clearpage
\section{flat and no very flat metrics}
\begin{colbox}
	\begin{enumerate}
		\item How many degrees of freedom does the metric have in \(d\) dimensions? How many degrees of freedom does a coordinate transofmation have in \(d\) dimensions?
		\item We'll look at an infinite cylinder of radius \(R\). We'll mark its surface with \(S\). \(S\) is a two dimensional space and can be described by a shape in three dimensional euclidian space:
		\begin{equation}
			x^2+y^2=R^2
		\end{equation}
		Move to cylindrical coordinates and find the metric for \(S\). Find a coordinate system where the metric is euclidian, i.e \(g\indices{_i_j} = \delta_{ij}\)
		\item Calculate the metric on the surface of a sphere with radius \(R\).
	\end{enumerate}
\end{colbox}
\subsection{degrees of freedom}
The metric has \(\frac{d\ins{d+1}}{2}\) degrees of freedom in \(d\) dimensions. This is because the metric is a symmetric tensor. The coordinate transformation has \(d\) degrees of freedom because there are \(d\) new coordinate functions.
\subsection{Cylinder}
We'll move to cylindrical corodinate such that
\begin{align}
	x = R\cos\theta && y=R\sin\theta && z=z
\end{align}
Therefore
\begin{align}
	\dd x = -R\sin\theta\dd\theta && \dd y = R\cos\theta\dd\theta && \dd z = \dd z
\end{align}
giving us
\begin{align}
	\dd l^2 = R^2\dd\theta^2 + \dd z^2
\end{align}
so the metric is
\begin{align}
	g\indices{_{ij}} = \pmat{R^2&0\\0&1}
\end{align}
if we move to a coordinate system such that
\begin{align}
	u=R\theta
\end{align}
we get
\begin{align}
	\dd l^2 = \dd u^2 + \dd z^2
\end{align}
or in other words
\begin{equation}
	g\indices{_{ij}} = \delta\indices{_{ij}}
\end{equation}
\subsection{surface of sphere}
In spherical coordinates the surface is given by
\begin{align}
	x=R\sin\theta\cos\varphi && y=R\sin\theta\sin\varphi && z=R\cos\theta
\end{align}
the line element in spherical coordinates is given by
\begin{equation}
	\dd l^2 = \dd r^2 + r^2\dd\theta^2 + r^2\sin^2\theta\dd\varphi^2
\end{equation}
Since \(r=R\) is constant \(\dd r=0\) so
\begin{equation}
	\dd l^2 = R^2\dd\theta^2 + R^2\sin^2\theta\dd\varphi^2
\end{equation}
so the metric is
\begin{equation}
	g\indices{_{ij}} = \pmat{R^2&0\\0&R^2\sin^2\theta}
\end{equation}